\documentclass[11pt]{article}
\usepackage[T1]{fontenc}
\usepackage{lmodern}
\usepackage[letterpaper,margin=0.75in]{geometry}
\usepackage{amsmath,amssymb,microtype,xcolor,hyperref}
\hypersetup{colorlinks=true,urlcolor=blue!45!black,linkcolor=blue!45!black,
 pdftitle={OpenAI QCD: assessment and renormalization prescription},
 pdfauthor={Review prepared with Codex},
 pdfsubject={Results, verification, mass-prescription discrepancy and unfinished checks}}
\setlength{\parindent}{0pt}
\setlength{\parskip}{5pt}
\pagestyle{plain}
\clubpenalty=10000
\widowpenalty=10000
\newcommand{\reviewdoc}[2]{\href{https://github.com/yang-mills-group/QCD/blob/main/#1}{#2}}

% Page 1: assessment. Page 2: renormalization prescription and comparison.
% Supporting links require repository access.
\begin{document}
\begin{center}
{\Large\bfseries OpenAI massive QCD: assessment}\par
\vspace{3pt}
25 November 2026
\end{center}
\vspace{-4pt}

\textbf{1. Results and relation to pure Yang--Mills.}
The Lean theorem extends the framework used in the November 4 pure
Yang--Mills submission to four-dimensional SU(3) with six fundamental Dirac
flavors. It uses nearest-neighbor Wilson fermions ($r=1$) and retains the
full quark determinants, so the fermions contribute to the gauge-field
weight. The development shares substantial lattice and analytic machinery
with the pure-gauge construction. It has its own final theorem and does
not directly invoke the earlier final existence or perturbative-matching
theorems; their conclusions do not automatically carry over to QCD.

The theorem supplies one common construction with continuum correlations,
reflection positivity, six positive distinct scheme masses, current and
stress-tensor identities, and a Wightman reconstruction. The mass and Ward
assertions survive selection of a vacuum sector. The source also exhibits
the one-loop flavor dependence $B_0=11-2N_f/3$ and fermion screening.
The result's potential significance is a continuum existence construction
with dynamical massive quarks. Establishing that its fields and parameters
have the conventional QCD interpretation requires further checks.

\textbf{2. Checks supplied by OpenAI.}
All 7,594 Lean modules compiled. Source hashes and comparison with frozen
theorem statements and definitions passed. Comparator checks that the
solution proves the preserved target, including its definitions; Nanoda
and Lean's kernel checked the proof with the allowed standard axioms.
We found no separate supplied numerical physics test suite.
These checks establish formal correctness of the specified statements;
they do not by themselves identify the defined masses and coupling with
experimentally normalized quantities.

\textbf{3. Our additional checks.}
We checked limits of actual lattice correlations at separated points,
properties retained in a selected vacuum, the exact mass fit and its
stability, bounds on density rescalings, flavor dependence, and a
conditional test of stress-tensor uniqueness. All 37 additional Comparator
targets passed both kernels. Independent finite-lattice diagnostics include
41 Wilson-operator checks and 44 exact mass-fit checks. Deliberately altered
inputs test whether the verification workflow detects errors.

Ten later Lean checks address Wilson contact identities and normalization;
controls with a wrong sign or an omitted Wilson remainder are rejected.
These ten have not had a Comparator run. Some added proofs derive
consequences of the submitted existence theorem; they do not provide a
second construction. All runs used macOS without the supplied Linux isolation.

\textbf{4. Assessment and remaining problems.}
The completed checks support the formal result and make several of its
consequences explicit. The main unresolved physical issue is the relation
between the prescribed operator normalization and conventional flavor
currents, including contact terms. The six masses are fitted outputs in
reference units. Neither arbitrary prescribed quark masses nor a conversion
to a conventional coupling, $\Lambda_{\rm QCD}$ or GeV has been established.
The formal all-pairs mass fit also differs in general from the six-pair rule
in the informal account. Our tests distinguish the rules on synthetic
readings; they do not show a discrepancy in the actual constructed theory.
The prescriptions and this distinction are explained on page~2.

The comparison with a free gauge field plus six free Dirac fields remains
unfinished: our conditional stress-tensor test does not complete that
comparison. Physical two-loop matching is deferred. Lean verification does
not settle the separate informal review's unresolved analytic estimates.
No particle spectrum, mass gap or confinement result is established.

{\small\textbf{Records:}
\reviewdoc{docs/qcd-verification-assessment.pdf}{Detailed assessment};
\reviewdoc{docs/evidence.md}{verification records};
\reviewdoc{docs/mass-current-normalization.md}{mass/current audit};
\reviewdoc{docs/informal-part-VI-review-r02.pdf}{informal review}.
Prepared with Codex. Links require repository access.}

\newpage
\begin{center}
{\Large\bfseries Renormalization prescription}\par
\vspace{3pt}
Operator units, fitted masses, coupling and scale
\end{center}
\vspace{-4pt}

\textbf{1. First normalize the operators.}
The axial current $A^\mu_{ij}$ and pseudoscalar density $P_{ij}$ correspond to
$\bar\psi_i\gamma^\mu\gamma_5\psi_j$ and $\bar\psi_i\gamma_5\psi_j$.
Integrate out short-distance variables down to the reference length~1.
Let $\rho$ be the resulting blocked density (the effective integration
weight), and $\rho_O$ the density obtained with an operator $O$ inserted.
Form $R_O=\rho^{-1}\rho_O$ near trivial gauge links, where the inverse
exists, before integrating the remaining gauge and Grassmann variables.

Choose the operator-mixing coefficients so that $R_O$, evaluated on the
prescribed smooth gauge and fermion test profiles, has the Taylor
coefficients of the corresponding \emph{unit-coefficient} bilinear through
the prescribed order. The other coefficients at those orders must also
match. This fixes the reference normalization of $A$ and $P$ in the chosen
scheme. Smooth gauge connections here supply test links; this is not an
interpolation of every fluctuating lattice gauge configuration.
Agreement with conventional flavor-current normalization, including
contact terms, remains unverified.

\textbf{2. Fit masses, then rescale the densities.}
With these initial units fixed, the theorem asserts
$\partial_\mu A^\mu_{ij}=b_{ij}P_{ij}$ for $i\ne j$. Let
$u=(u_1,\ldots,u_6)$ be six \emph{trial masses}; $m$ is the unique best fit
to all 30 ordered off-diagonal pairs:
\[
 m=\underset{u\in\mathbb R^6}{\operatorname{argmin}}
       \sum_{i\ne j}(b_{ij}-u_i-u_j)^2.
\]
The fitted masses are close to positive, separated readings from the
actual blocked density. They are outputs in reference units, distinct from
the bare Wilson mass parameters; no numerical quark or particle masses are
specified. Now define an \emph{additional}, pair-dependent rescaling,
\[
 P^R_{ij}:=\frac{b_{ij}}{m_i+m_j}P_{ij},
 \qquad
 \partial_\mu A^\mu_{ij}=(m_i+m_j)P^R_{ij}.
\]
Thus $P^R$ is the already calibrated density $P$ multiplied by the displayed
factor; $A$ is unchanged. The final identity follows algebraically from the
unrescaled one. Scalar densities are similarly rescaled to put the vector
identity in terms of $m_i-m_j$.

\textbf{3. Coupling and scale.}
The reference length is $L_0=1$, with $\beta(0)=\beta_{\rm ref}$;
blocking ratios specify finer spacings. The inverse-coupling coordinate
$\beta$ is required to satisfy $|\kappa-\beta|\leq\varepsilon$, where
$\kappa$ is a normalized second-derivative reading near trivial gauge links
of $-\log\operatorname{Re}\rho_0$. Here $\rho_0$ is the scalar Grassmann
coefficient of the blocked density. Thus the prescription is a tolerance
band, not exact equality to that reading. Existence holds for sufficiently
large $\beta_{\rm ref}$ after choosing $\varepsilon>0$. No conversion to an
observable coupling, a conventional $\Lambda_{\rm QCD}$, or a scale in GeV
has been established by this review.

\textbf{4. Formal versus informal masses.}
The informal account, as recorded in the Part VI review (p.~8), uses only
$\{12,13,23,14,15,16\}$:
\[
 m_1=(b_{12}+b_{13}-b_{23})/2,
 \qquad m_j=b_{1j}-m_1\quad(j=2,\ldots,6).
\]
It then rescales the remaining densities. Both rules agree if every original
$b_{ij}$ is an exact mass sum. Our synthetic perturbation of $b_{56}=b_{65}$
leaves the six-pair answer unchanged but changes the all-pairs fit.
The prescriptions therefore differ in general; disagreement on the actual
constructed correlations has not been shown.

{\small\textbf{Definitions:}
\reviewdoc{provenance/supplied-comparator/frozen/reference/YM/Massive/Target/Calibration.lean}{Operator calibration};
\reviewdoc{provenance/supplied-comparator/frozen/reference/YM/Massive/Target/Physics.lean}{mass and current prescription};
\reviewdoc{provenance/supplied-comparator/frozen/reference/YM/Massive/Target/Tower.lean}{coupling and scale}.
Supporting discussion is in the records linked on page~1.}
\end{document}
